\documentclass{article}

\usepackage{neurips_2026}
\usepackage[utf8]{inputenc}
\usepackage[T1]{fontenc}
\usepackage{hyperref}
\usepackage{url}
\usepackage{booktabs}
\usepackage{amsfonts}
\usepackage{nicefrac}
\usepackage{microtype}
\usepackage{xcolor}
\usepackage{graphicx}

\title{Who Do They Cite? Gendered Citation Patterns in French Legislative Manifestos (1973--1993)}

\author{%
  Anahí Reyes \\
  M2 Machine Learning for NLP \\
  École Polytechnique \\
  \texttt{anahireyes08@gmail.com}
}

\begin{document}

\maketitle

\begin{abstract}
Political manifestos are a window into how candidates construct their political identity
and position themselves within a broader discourse. This paper asks whether male and
female candidates in French legislative elections cite political figures differently,
with a focus on the gender of cited persons. Using a corpus of 20,331 manifestos from
five elections (1973--1993), we build an NLP pipeline that extracts person citations
via a CamemBERT-NER model and resolves their gender through Wikidata. We find a
large and statistically significant gap: female candidates cite women in 44\% of their
known-gender citations, compared to 5\% for male candidates
(Mann-Whitney $U$, $p < 0.001$). This gap persists across all five election years.
Beyond the quantitative difference, the \emph{identity} of cited figures reveals a
qualitative distinction: male candidates cite women as electoral actors, while female
candidates draw on feminist intellectuals and activists absent from male manifestos.
We interpret this as evidence that female candidates actively construct a female
political tradition in their discourse, compensating for its near-absence in the
institutional political landscape.
\end{abstract}


%────────────────────────────────────────────────────────────
\section{Introduction}
%────────────────────────────────────────────────────────────

Political representation is not only about who holds office, but about who is made
visible in political discourse. When a candidate writes a manifesto, the figures they
choose to cite — by name — signal their political genealogy, their ideological
commitments, and their vision of who belongs in the political sphere.

France between 1973 and 1993 is a particularly revealing setting. Female candidacies
were rare (around 10\% of the corpus) and institutionally marginalised, yet this period
saw the emergence of prominent feminist figures — Gisèle Halimi, Arlette Laguiller,
Huguette Bouchardeau — who became central reference points in left-wing politics.
Did female candidates reference these figures more than their male counterparts?
And more broadly, do male and female candidates inhabit different discursive worlds
when it comes to whom they name in their manifestos?

We address these questions by combining named entity recognition (NER) over a large
corpus of digitised manifestos with a structured knowledge base (Wikidata) to
assign gender to cited persons. Our contributions are:

\begin{itemize}
  \item A replicable NLP pipeline for person citation extraction and gender tagging
        at scale, applicable to any corpus of political texts.
  \item Quantitative evidence of a large, stable gender gap in citation behaviour
        across five French legislative elections.
  \item A qualitative analysis of which figures are cited, revealing that the gap is
        not only about quantity but about the \emph{type} of female figures invoked.
\end{itemize}


%────────────────────────────────────────────────────────────
\section{Related Work}
%────────────────────────────────────────────────────────────

Gender gaps in political representation have been studied extensively in political
science \cite{norris1997passages, celis2008}. Research on political discourse has
shown that male and female politicians differ in rhetorical style, issue emphasis,
and self-presentation \cite{lawless2004}. However, citation behaviour in political
manifestos — who candidates name and why — has received less attention as a lens on
gendered discourse.

The role model effect in political participation suggests that the visibility of
female political figures increases other women's sense of political efficacy
\cite{campbell2010}. Our work operationalises this idea at the text level: which
figures do female candidates make visible in their manifestos?

On the NLP side, NER for French political texts has been explored using
CamemBERT-based models \cite{martin2020camembert}, and knowledge base linking for
gender attribution via Wikidata is an established approach \cite{wagner2016}. We
combine these tools in a pipeline designed for historical French political documents,
which present particular challenges: OCR noise, non-standard formatting, and
multilingual content (some Alsatian manifestos are in German).


%────────────────────────────────────────────────────────────
\section{Data}
%────────────────────────────────────────────────────────────

\subsection{Archelec Corpus}

We use the Archelec corpus, a digitised collection of French legislative election
manifestos (\emph{professions de foi}) from the French National Archives. We focus
on five legislative elections: 1973, 1978, 1981, 1988, and 1993 — the years for
which digitised text is available. The corpus comprises 20,331 manifestos after
joining with candidate metadata (gender, party, department).

\begin{table}[h]
  \caption{Corpus composition by election year and candidate gender.}
  \label{tab:corpus}
  \centering
  \begin{tabular}{lrrr}
    \toprule
    Year & Total manifestos & Female candidates & \% Female \\
    \midrule
    1973 & 3,921 & -- & -- \\
    1978 & 5,030 & -- & -- \\
    1981 & 3,182 & -- & -- \\
    1988 & 3,628 & -- & -- \\
    1993 & 5,936 & -- & -- \\
    \midrule
    Total & 20,331 & 2,160 & 10.0\% \\
    \bottomrule
  \end{tabular}
\end{table}

Female candidates represent 10\% of the corpus, consistent with the historical
record of female candidacy rates in France during this period, which peaked at
13.6\% in 1978 before declining. Digitisation coverage is virtually identical
across genders ($\approx$100\%), ruling out a sampling bias.

\subsection{Text Quality}

OCR noise — measured as the ratio of non-alphabetic, non-space characters — is
low and stable across years (mean $\approx$0.045). Language detection identifies
97.8\% of manifestos as French; 1.9\% are in English (likely misclassified) and
0.16\% in German (Alsace-Moselle). Non-French manifestos are retained but NER
performance may be lower on these.

The median manifesto is 963 tokens; 88\% exceed the 512-token limit of
CamemBERT-NER, meaning NER covers only the first $\approx$350 words of most
texts. We accept this constraint, noting that candidates typically name key
political figures early in their manifestos.


%────────────────────────────────────────────────────────────
\section{Methods}
%────────────────────────────────────────────────────────────

\subsection{Named Entity Recognition}

We extract person citations using \texttt{Jean-Baptiste/camembert-ner}, a
CamemBERT model fine-tuned on French NER datasets (WikiNER and the French
TreeBank). We apply it with \texttt{aggregation\_strategy=``simple''} to merge
subword tokens into full entity spans, retaining only \texttt{PER} entities with
confidence $> 0.7$.

\paragraph{NER evaluation.}
We manually annotated a stratified sample of 50 manifestos (25 female, 25 male),
generating 154 model predictions. Annotators marked each prediction as correct
(TP) or erroneous (FP) and noted missed names (FN). Results:

\begin{center}
  Precision = 0.857 \quad Recall = 0.978 \quad F1 = 0.913
\end{center}

Precision is lower for female-candidate manifestos (0.812 vs.\ 0.913 for male),
likely because female manifestos contain more local and non-public figures that
the model tags with lower accuracy.

\subsection{Citation Cleaning}

Raw NER output contains 85,827 citations (38,460 unique names). We apply four
cleaning steps:

\begin{enumerate}
  \item \textbf{Noise removal} (--4,910): honorifics and titles misclassified as
        persons (\emph{Monsieur}, \emph{Président de la République}, etc.)
  \item \textbf{Single-token removal} (--10,470): one-word spans are too
        ambiguous to identify reliably.
  \item \textbf{Truncated-name removal} (--441): names ending in a single
        character, e.g.\ ``Giscard D'', are NER/OCR artefacts.
  \item \textbf{Hapax removal} (--19,590): names appearing only once across the
        entire corpus are likely noise or hyper-local figures.
\end{enumerate}

After cleaning: 50,416 citations, 10,139 unique names.

\subsection{Gender Tagging via Wikidata}

For each of the 10,139 unique names, we query the Wikidata MediaWiki API in two
steps: (1) \texttt{wbsearchentities} to retrieve candidate Q-IDs, and (2)
\texttt{wbgetentities} to verify that the entity is a human
(\texttt{wdt:P31 wd:Q5}) and retrieve their sex or gender property
(\texttt{wdt:P21}). A 0.3-second delay between requests was added to comply with
Wikidata rate limits.

We resolved gender for 5,146 / 10,324 queried names (49.8\%), yielding:

\begin{center}
  \begin{tabular}{lr}
    \toprule
    Gender & Citations \\
    \midrule
    Male    & 30,995 (61.5\%) \\
    Unknown & 17,073 (33.9\%) \\
    Female  &  2,337 (4.6\%)  \\
    Other   &     11 (0.0\%)  \\
    \bottomrule
  \end{tabular}
\end{center}

\paragraph{Limitation.}
The 33.9\% unknown rate reflects Wikidata's documented gender gap: women,
particularly local politicians and minor party figures, are systematically less
likely to have Wikidata entries. All subsequent analyses are computed on the
resolved subset only (66.1\% of citations) and should be interpreted as
conservative lower bounds, since unknown citations likely contain a
disproportionate share of women.


%────────────────────────────────────────────────────────────
\section{Results}
%────────────────────────────────────────────────────────────

\subsection{Core Finding: Female Candidates Cite Women Significantly More}

For each candidate, we compute the share of their known-gender citations going
to female figures. We compare the distributions for female and male candidates
using a one-sided Mann-Whitney $U$ test (non-parametric, robust to the heavily
skewed distributions observed).

\begin{center}
  \begin{tabular}{lrr}
    \toprule
    & Female candidates & Male candidates \\
    \midrule
    $N$ candidates & 1,320 & 12,262 \\
    Mean \% female cited & 43.99\% & 4.69\% \\
    \midrule
    \multicolumn{3}{l}{Mann-Whitney $U = 12{,}463{,}618$, $p < 0.001$} \\
    \bottomrule
  \end{tabular}
\end{center}

Female candidates cite women nearly nine times more often than male candidates
do, and this difference is highly significant. The gap of $\approx$39 percentage
points is striking given that female candidates and male candidates navigate the
same mainstream political landscape — both groups cite Mitterrand, Rocard, and
Le Pen as their dominant references.

\subsection{Temporal Stability}

The gap persists across all five election years examined, with female candidate
rates ranging from 28\% (1988) to 58\% (1978) and male candidate rates ranging
from 2\% (1973) to 7\% (1978). The gap never closes, ranging from 26 to 51
percentage points depending on the year.

A complementary analysis of \emph{unique} cited figures (deduplicating names
per group-year) reveals a steadily increasing repertoire of female figures for
female candidates (40\% in 1973 to 48\% in 1988), while the unique-figure rate
for male candidates remains flat at 2--5\%. This divergence indicates that the
volatility in citation-event rates reflects concentration effects (a few women
cited repeatedly) rather than a genuine decline in female citation. Female
candidates were steadily broadening the diversity of women they cited over time.

\subsection{Who Is Cited? Qualitative Pattern}

Both groups share the same dominant references: François Mitterrand
(20.5\% of female manifestos, 15.7\% of male), Michel Rocard, and Le Pen.
The mainstream political landscape looks the same for both genders. The
differences emerge at the margins but are qualitatively significant.

\paragraph{Figures distinctive to female candidates.}
Arlette Laguiller (Lutte Ouvrière, far-left) appears in 16.2\% of female
manifestos vs.\ 2.7\% of male — a sixfold difference. Gisèle Halimi (feminist
lawyer), Simone de Beauvoir (feminist philosopher), and Huguette Bouchardeau
(PSU, left-feminist politician) appear in 2--3\% of female manifestos but are
absent from the male top 20. Humanist intellectuals Jacques Monod and Jean
Rostand (biologist and philosopher-scientist) also appear exclusively in female
candidates' top references.

\paragraph{Figures distinctive to male candidates.}
Male-exclusive references are exclusively orthodox party politicians: Pierre
Mauroy, Raymond Barre, Georges Marchais, Antoine Pinay, Robert Fabre, Michel
Jobert, Jean Lecanuet — all men, all institutional figures from across the
left-right spectrum.

\paragraph{Interpretation.}
Male candidates cite women only as conventional electoral actors (a rival
presidential candidate, a minister who passed a famous law). Female candidates
cite women as feminist intellectuals, activists, and advocates — figures who
represent not who holds power, but who has theorised or fought against its
gendered distribution. The citation gap is therefore not only quantitative but
qualitative: female candidates use citations to construct a feminist political
identity, drawing on a non-electoral intellectual tradition precisely because
the electoral space offers so few female role models.

This distinction is an inference from the political profiles of the cited
figures, not from textual analysis of the surrounding sentences. Confirming the
rhetorical function of these citations would require close reading of manifesto
passages — a natural extension of this work.

\subsection{Party-Level Variation}

Controlling for party family reveals a left-right gradient: left-wing parties
(PCF, PS, PSU) exhibit higher female citation rates than centrist and right-wing
parties (Centre/MRP, RPR, UDF). The feminist party \emph{Choisir la cause des
femmes} (founded by Gisèle Halimi) stands as an extreme case. This pattern is
consistent with the broader literature on left parties' greater formal commitment
to gender equality in political representation during this period.


%────────────────────────────────────────────────────────────
\section{Discussion}
%────────────────────────────────────────────────────────────

Our results support a structural rather than an individual explanation for the
low overall female citation rate. Male candidates are not necessarily hostile to
citing women — they simply operate in a political world where women had not yet
reached the level of national prominence required to become a natural reference
point in electoral discourse. The pool of citable female figures was structurally
scarce.

Female candidates respond to this scarcity by reaching beyond the electoral
sphere. By citing Halimi and Beauvoir alongside Laguiller, they are actively
populating the political landscape with female reference points — building a
tradition where one was not institutionally provided. This is consistent with the
role model theory of political participation \cite{campbell2010}: female
candidates may be using their manifestos to signal to female readers that there
is a lineage, a set of women who fought and thought about power, and that
political involvement is a legitimate space for women.

\paragraph{Limitations.}
Three limitations deserve emphasis. First, the 33.9\% unknown gender rate
introduces a conservative bias: if unknown citations skew female (as
Wikidata's gender gap suggests), the true citation gap is larger than reported.
Second, NER covers only the first $\approx$350 words of most manifestos, missing
citations in later paragraphs. Third, citation counts are not citation context:
we cannot determine from this analysis whether a name is cited approvingly,
critically, or neutrally. Fourth, the study period ends in 1993; it would be
valuable to extend the analysis to more recent elections, where female candidacy
rates have grown substantially.

\paragraph{Future work.}
A name-based gender classifier (e.g., \texttt{gender\_guesser}) applied to
unresolved citations would provide corrected estimates and directly quantify
the Wikidata bias. Extending the corpus to post-1993 elections would test
whether the citation gap narrows as more women enter the political mainstream.
Finally, sentence-level analysis of citation context would allow us to
distinguish symbolic citations (invoked to construct identity) from adversarial
ones (cited as rivals to criticise).


%────────────────────────────────────────────────────────────
\section{Conclusion}
%────────────────────────────────────────────────────────────

We have shown that female and male candidates in French legislative elections
(1973--1993) inhabit the same political landscape in their manifestos — the same
dominant male figures anchor both groups' discourse — but differ sharply in the
female figures they invoke. Female candidates cite women nearly nine times more
often, and the specific figures they choose reveal a deliberate effort to
construct a feminist political genealogy. This gap is stable across two decades
and robust to the party composition of the corpus.

The finding reframes a simple ``gender gap in citations'' as a deeper point about
political discourse: female candidates are doing something qualitatively distinct
with their manifestos, drawing on feminist thought and activism to populate a
political world that had structurally excluded women. The manifestos are not just
electoral documents; they are also, for female candidates, acts of political
identity-making.


\begin{ack}
This work was conducted as a final project for the M2 Machine Learning for NLP
course at École Polytechnique. The Archelec corpus is provided by Sciences Po /
CEVIPOF. Gender data is retrieved from Wikidata under CC0.
\end{ack}

\section*{References}

{
\small

[1] Campbell, R., \& Wolbrecht, C.\ (2006). See Jane run: Women politicians as
role models for adolescents. \emph{Journal of Politics}, 68(2), 233--247.

[2] Celis, K., et al.\ (2008). Rethinking women's substantive representation.
\emph{Representation}, 44(2), 99--110.

[3] Lawless, J.L.\ (2004). Politics of presence? Congresswomen and symbolic
representation. \emph{Political Research Quarterly}, 57(1), 81--99.

[4] Martin, L., et al.\ (2020). CamemBERT: a tasty French language model.
\emph{Proceedings of ACL 2020}.

[5] Norris, P., \& Lovenduski, J.\ (1995). \emph{Political Recruitment: Gender,
Race and Class in the British Parliament}. Cambridge University Press.

[6] Wagner, C., et al.\ (2016). It's a man's Wikipedia? Assessing gender
inequality in an online encyclopedia. \emph{Proceedings of ICWSM 2016}.

}

\appendix

\section{Data Attrition Summary}

\begin{verbatim}
Raw NER output:           85,827 citations  (38,460 unique names)

── Cleaning (notebook 02) ──────────────────────────────────
  Noise removed:          -4,910   (titles, honorifics)
  Single-word removed:   -10,470   (single tokens)
  Truncated names:          -441   (e.g. "Giscard D")
  Rare (< 2 times):      -19,590   (appear only once)
────────────────────────────────────────────────────────────
Clean citations:           50,416  (10,139 unique names)

── Wikidata gender tagging (notebook 03) ───────────────────
  Names queried:           10,324
  Found on Wikidata:        5,146 / 10,324  (49.8%)

  Gender assigned:
    male:                  30,995 citations
    female:                 2,337 citations
    non-binary/other:          11 citations
    unknown:               17,073 citations
────────────────────────────────────────────────────────────
Final coverage:    66.1% of citations have known gender
\end{verbatim}

\newpage
\input{checklist.tex}

\end{document}
